\NSPage{133}%
\begin{EESegment}{EE-TGT-P133-001}
On the pulse, \NSi{NS-F-P133-001}{XE^2=X_{\mathrm p} e_b^2 f^2 e^{-2\lambda y}}. Change variables to
\NSi{NS-F-P133-002}{\zeta_b=\lambda y}. Equations~\eqref{eq:A.14}, \eqref{eq:A.15}, and
\eqref{eq:A.18} then give the exact formula
\end{EESegment}
\NSFormula{NS-F-P133-003}{display}{A.19}
\begin{equation}
 \frac{\lambda S(\infty)}{X_{\mathrm p} e_b^2 f^2}
 =\operatorname{Amp}^2K_b-\frac{1-e^{-26}}{4}+\mathcal E(\operatorname{Amp},\eta),
 \qquad
 K_b=\int_0^{13}e^{-2\zeta_b}R_0(\zeta_b)^2\,d\zeta_b,
 \tag{A.19}\label{eq:A.19}
\end{equation}
\begin{EESegment}{EE-TGT-P133-002}
where, on \NSi{NS-F-P133-004}{[.9,1.2]\times[-1,1]},
\end{EESegment}
\NSFormula{NS-F-P133-005}{display}{none}
\[
 \lVert\mathcal E\rVert_{C^1_\eta}
 \leq C_{\mathrm{pre}}\lambda(1+\log(1/\lambda)).
\]
\begin{EESegment}{EE-TGT-P133-003}
Every term in this estimate that comes from the outer profile \NSi{NS-F-P133-006}{E} was chosen without using
\NSi{NS-F-P133-007}{\operatorname{Amp}}. Its only remaining dependence on the amplitude comes from the affine end
bumps, and that dependence is exponentially small. Since \NSi{NS-F-P133-008}{R_0\leq\zeta_b},
\NSi{NS-F-P133-009}{K_b\leq\int_0^\infty\zeta_b^2e^{-2\zeta_b}\,d\zeta_b=1/4}. On
\NSi{NS-F-P133-010}{.02\leq\zeta_b\leq9}, the bound goes the other way,
\NSi{NS-F-P133-011}{R_0\geq\zeta_b-.02}, and integration gives
\NSi{NS-F-P133-012}{K_b>.20}. The main expression in \eqref{eq:A.19} is less than
\NSi{NS-F-P133-013}{-.047} at amplitude \NSi{NS-F-P133-014}{.9}, and greater than
\NSi{NS-F-P133-015}{.038} at amplitude \NSi{NS-F-P133-016}{1.2}. Its derivative with respect to the
amplitude is at least \NSi{NS-F-P133-017}{.36} throughout that bracket. It follows that, when
\NSi{NS-F-P133-018}{\lambda} is small, the full expression has exactly one root in the bracket. The
linearization at that root is nonsingular, so it can be inverted locally. This makes the root smooth, and
\end{EESegment}
\NSFormula{NS-F-P133-019}{display}{A.20}
\begin{equation}
 |\partial_\eta\operatorname{Amp}|
 \leq C_{\mathrm{pre}}\lambda(1+\log(1/\lambda)).
 \tag{A.20}\label{eq:A.20}
\end{equation}
\begin{EESegment}{EE-TGT-P133-004}
Both equalities in \eqref{eq:A.8} now hold. What remains is to check that these moment corrections keep the stress
inside its cone. The bounds below prescribe smallness for only finitely many orders of
\NSi{NS-F-P133-020}{\eta}-derivatives. If smallness bounds are also set for any other finite list of derivative
orders, differentiating the integrals and applying Lemma~\ref{lem:A.2} gives those bounds before the parameters are
fixed. Once the parameters have been chosen, every other fixed derivative order still has a finite bound.
\end{EESegment}

\begin{EESegment}{EE-TGT-P133-005}
\subsection{Fixing the pressure datum for the axis profile.}\label{sec:A.4}
\end{EESegment}
\begin{EESegment}{EE-TGT-P133-006}
The azimuthal profile \NSi{NS-F-P133-021}{E} now fixes the pressure
\NSi{NS-F-P133-022}{\Pi_0} on the axis by requiring the pressure to vanish at radial infinity, as in
\eqref{eq:4.25}. Proposition~\ref{prop:B.2} constructs the axis profile analytically. For that construction to work,
this pressure datum must extend holomorphically to a neighborhood of \NSi{NS-F-P133-023}{[-1,1]} and must obey the
sign and size bounds in Lemma~\ref{lem:A.5}. This subsection checks those facts before the axis is attached. It also
shows that the datum does not depend on the later choice of radial scale \NSi{NS-F-P133-024}{X_R}. The pressure
integrals here use the global logarithmic coordinate \NSi{NS-F-P133-025}{y=\log(X/X_R)} on the full reference profile.
\end{EESegment}

\begin{lemma}\label{lem:A.5}
\begin{EESegment}{EE-TGT-P133-007}
Let \NSi{NS-F-P133-026}{E_{\mathrm{id,sched}}} consist of the reference inner profile \eqref{eq:A.7}, followed by the
outer profile \NSi{NS-F-P133-027}{E} from Section~\ref{sec:A.2}. Leave out the azimuthal corrections that preserve
the total pressure increment, and leave every other transition interval unchanged. Then the function
\end{EESegment}
\NSFormula{NS-F-P133-028}{display}{A.21}
\begin{equation}
 \Pi_0(\eta)=-\frac12\int_{-\infty}^{\infty}
 E_{\mathrm{id,sched}}(y,\eta)^2\,dy
 \tag{A.21}\label{eq:A.21}
\end{equation}
\begin{EESegment}{EE-TGT-P133-008}
does not depend on \NSi{NS-F-P133-029}{X_R} and is analytic on a complex neighborhood of
\NSi{NS-F-P133-030}{[-1,1]}. It is even. The derivative \NSi{NS-F-P133-031}{\Pi_0'(\eta)} has the same sign as
\NSi{NS-F-P133-032}{\eta} whenever \NSi{NS-F-P133-033}{\eta\neq0}. It also satisfies
\end{EESegment}
\NSFormula{NS-F-P133-034}{display}{A.22}
\begin{equation}
 \Pi_0(\eta)\leq-\frac52P_*^2f(\eta)^2.
 \tag{A.22}\label{eq:A.22}
\end{equation}
\begin{EESegment}{EE-TGT-P133-009}
For the full reference radial profile, start with this datum and integrate
\NSi{NS-F-P133-035}{\Pi_X=E^2/(2X)}. The result is
\end{EESegment}
\NSFormula{NS-F-P133-036}{display}{A.23}
\begin{equation}
 \Pi(X,\eta)=-\frac12\int_{\log(X/X_R)}^\infty E(v,\eta)^2\,dv.
 \tag{A.23}\label{eq:A.23}
\end{equation}
\end{lemma}
\begin{proof}
\begin{EESegment}{EE-TGT-P133-010}
Without the angular bumps, every part of \NSi{NS-F-P133-037}{E} can be written as
\NSi{NS-F-P133-038}{c(y)f(\eta)^{\vartheta(y)}}. Here \NSi{NS-F-P133-039}{c} and the transition lengths do not depend
on \NSi{NS-F-P133-040}{\eta}, and \NSi{NS-F-P133-041}{0\leq\vartheta\leq1}. The reference inner profile, the axial
transition, the intermediate power-law interval, and the pulse all have
\NSi{NS-F-P133-042}{\vartheta=1}. The profile interpolation lowers this exponent to zero, and the exterior transition
and tail have \NSi{NS-F-P133-043}{\vartheta=0}. The next step is to choose one simply connected neighborhood of the
real interval in the complex plane, one connected region in which every closed loop can shrink to a point without
leaving the region. Once the poles and zeros are excluded, this property lets one logarithm be chosen consistently
throughout the region.
\end{EESegment}
\NSPage{134}%
\begin{EESegment}{EE-TGT-P134-001}
Choose it so that it avoids every pole and zero of \NSi{NS-F-P134-001}{f}, and fix an analytic logarithm there, meaning
a holomorphic function whose exponential is the original function. Then
\NSi{NS-F-P134-002}{f^\vartheta=\exp(\vartheta\log f)} is holomorphic, with a bound that is uniform for
\NSi{NS-F-P134-003}{0\leq\vartheta\leq1}. The squared integrand is bounded by
\NSi{NS-F-P134-004}{Ce^{y/5}} as \NSi{NS-F-P134-005}{y\to-\infty}, and by
\NSi{NS-F-P134-006}{Ce^{-(1+2h)y}} as \NSi{NS-F-P134-007}{y\to+\infty}. The interval between those two tails is
finite. Integrating the holomorphic functions with these uniform bounds proves that the pressure datum is analytic.
\end{EESegment}

\begin{EESegment}{EE-TGT-P134-002}
For real \NSi{NS-F-P134-008}{\eta}, differentiating each factor gives
\NSi{NS-F-P134-009}{\partial_\eta f^{2\vartheta}=-2\vartheta J_0'f^{2\vartheta}}. It follows that every term in
\NSi{NS-F-P134-010}{\Pi_0'} has the same sign as \NSi{NS-F-P134-011}{\eta}. The reference inner profile makes the
sign strict. That profile contributes exactly \NSi{NS-F-P134-012}{-(5/2)P_*^2f^2}, which proves
\eqref{eq:A.22}. Nothing in the construction in the logarithmic radial coordinate contains
\NSi{NS-F-P134-013}{X_R}, so the pressure datum does not depend on that scale. Finally, \eqref{eq:A.11} says that putting the
angular bumps back changes the total pressure increment by zero. So the full increment is
\NSi{NS-F-P134-014}{-\Pi_0}, which proves \eqref{eq:A.23}.
\end{EESegment}
\end{proof}

\begin{EESegment}{EE-TGT-P134-003}
The axis construction can now use this pressure datum. Every later change to the profile must keep the total pressure
increment unchanged so that the datum stays fixed. Matching all five radial integrals also keeps the radial velocity and
the functions \NSi{NS-F-P134-015}{Q_s,N_s} unchanged past the correction interval. Suppose an edit has zero total
integral of \NSi{NS-F-P134-016}{(E_{\mathrm{new}}^2-E_{\mathrm{old}}^2)/(2X)}. Then the pressure found by integrating
forward is unchanged before and after the edit's support, although it may change inside that support. A correction to the
strict right of \NSi{NS-F-P134-017}{X} can be left out of the backward integral \eqref{eq:A.23} when it
preserves the total pressure increment. Lemma~\ref{lem:4.4} gives a stronger result when all five radial moment functions
agree at a radius where the two fields meet again. From that radius onward, the radial velocity and the functions
\NSi{NS-F-P134-018}{Q_s,N_s} are restored on the whole later interval, along with their parameter derivatives. The axis
construction uses the analytic function \NSi{NS-F-P134-019}{\Pi_0}, which is already fixed, and integrates pressure
forward. Later edits of \NSi{NS-F-P134-020}{E} may be smooth without being analytic, but they cannot change that datum.
Their total pressure increment is either zero or is restored by the correction that matches all five radial moments. The
heat replacement changes the exterior profile while preserving the required convergent differences of moments.
\end{EESegment}

\begin{EESegment}{EE-TGT-P134-004}
\subsection{Uniform admissible stress-cone inequalities on the outer interval.}\label{sec:A.5}
\end{EESegment}
\begin{EESegment}{EE-TGT-P134-005}
The moment identities and the pressure datum on the axis are now fixed. It remains to check the stress cone inequalities
on the outer intervals named in Proposition~\ref{prop:A.4}. Before choosing the large radial scale
\NSi{NS-F-P134-021}{X_R}, consider the ratios in the sufficient cone test from Lemma~\ref{lem:4.5}. The final choice of
that scale will then make the inequalities strict, with the same margin throughout the required finite interval.
\end{EESegment}

\begin{EESegment}{EE-TGT-P134-006}
Write \NSi{NS-F-P134-022}{w=N_s/(EQ_s)} wherever \NSi{NS-F-P134-023}{Q_s>0}. The sufficient test from
Lemma~\ref{lem:4.5} is
\end{EESegment}
\NSFormula{NS-F-P134-024}{display}{A.24}
\begin{equation}
 a-b_sw>0,\qquad 2b_sw+b_s^2/a+(a-2)w^2<2.
 \tag{A.24}\label{eq:A.24}
\end{equation}
\begin{EESegment}{EE-TGT-P134-007}
The proof will give uniform positive lower bounds for \NSi{NS-F-P134-025}{a-b_sw} and
\NSi{NS-F-P134-026}{2-2b_sw-b_s^2/a-(a-2)w^2}. Throughout this check,
\NSi{NS-F-P134-027}{a,b_s,w} stay in fixed compact ranges and \NSi{NS-F-P134-028}{Q_s>0}. Afterward, increasing
\NSi{NS-F-P134-029}{X_R} makes \NSi{NS-F-P134-030}{p_{s,1}=XQ_s/L} large enough throughout the chosen finite
logarithmic interval. At points where \NSi{NS-F-P134-031}{v_s\leq2}, the remaining requirement is
\NSi{NS-F-P134-032}{P_{\mathrm c}=p_{s,1}(1-b_sw/a)>2}.
\end{EESegment}

\begin{EESegment}{EE-TGT-P134-008}
\emph{The reference inner profile and axial transition.} On \eqref{eq:A.7}, putting the profile directly into the
equation for \NSi{NS-F-P134-033}{Q_s} gives
\end{EESegment}
\NSFormula{NS-F-P134-034}{display}{A.25}
\begin{equation}
 Q_s=\frac{(3/5)(4L-1)-h(1-8\eta^2)+(D+4d)\eta J_0'}{8/5}\geq c>0.
 \tag{A.25}\label{eq:A.25}
\end{equation}
\begin{EESegment}{EE-TGT-P134-009}
There, \NSi{NS-F-P134-035}{b_s=0} and \NSi{NS-F-P134-036}{a=4/5}. Split the source in the equation for
\NSi{NS-F-P134-037}{N_s} into its geometric part and its pressure part. At the beginning of the first transition, the
geometric part is \NSi{NS-F-P134-038}{O(|\eta|)}, and the pressure part is
\NSi{NS-F-P134-039}{O(P_*^2|\eta|)}. One parameter
derivative removes the factor \NSi{NS-F-P134-040}{|\eta|}. These bounds come from the constant ideal geometric source
and the bounded pressure source.
\end{EESegment}
\NSPage{135}%
\begin{EESegment}{EE-TGT-P135-001}
Both sources are integrated against the kernel in the integral representation of \NSi{NS-F-P135-001}{N_s}. The
pressure and ideal shapes are even, so both sources vanish at \NSi{NS-F-P135-002}{\eta=0}.
\end{EESegment}

\begin{EESegment}{EE-TGT-P135-002}
During the first transition, \NSi{NS-F-P135-003}{U=4\eta},
\NSi{NS-F-P135-004}{W=1-4L<0}, and \NSi{NS-F-P135-005}{l\geq0}, so the angular source is at least
\NSi{NS-F-P135-006}{-h}. Since the transition has finite length, \NSi{NS-F-P135-007}{Q_s} keeps a positive lower
bound. On this interval, \NSi{NS-F-P135-008}{b_s=0} and \NSi{NS-F-P135-009}{a\leq2}. In the axial transition, the
angular source \NSi{NS-F-P135-010}{S_q} is
\end{EESegment}
\NSFormula{NS-F-P135-011}{display}{A.26}
\begin{equation}
 S_q=-h(1-2k\eta^2)+(D+dk)\eta J_0',\qquad
 Q_s\geq c(\eta^2+e^{-y})\quad(0\leq y\leq T_d).
 \tag{A.26}\label{eq:A.26}
\end{equation}
\begin{EESegment}{EE-TGT-P135-003}
Solving the linear equation by variation of constants---multiplying by its integrating factor and integrating---gives
the second bound. The facts used here are
\NSi{NS-F-P135-012}{\eta J_0'\asymp\eta^2}, \NSi{NS-F-P135-013}{D+dk\geq D}, positivity of
\NSi{NS-F-P135-014}{Q_s} at the inner endpoint, and \NSi{NS-F-P135-015}{h\ll e^{-T_d}}.
\end{EESegment}

\begin{EESegment}{EE-TGT-P135-004}
On this interval the geometric axial source is \NSi{NS-F-P135-016}{O(|\eta|)}. Equation~\eqref{eq:A.23} makes its
pressure source \NSi{NS-F-P135-017}{O(E^2|\eta|)}. The reason is that the unmodified profile at larger radii decays at
least as fast as a constant times \NSi{NS-F-P135-018}{e^{-\Delta y/2}} and satisfies
\NSi{NS-F-P135-019}{|\partial_\eta\log E|\leq|J_0'|}. The later correction to the azimuthal profile can be left out
because it preserves the total pressure increment. Solve \NSi{NS-F-P135-020}{N_s'+N_s=S_n} and use
\NSi{NS-F-P135-021}{E^2=P_1^2f^2e^{-y}}. This gives
\end{EESegment}
\NSFormula{NS-F-P135-022}{display}{none}
\[
 \frac{|N_s|}{E^2}\leq C|\eta|
 \left(\frac{e^{T_d}}{P_*^2}+1+y\right)\leq C|\eta|(1+y).
\]
\begin{EESegment}{EE-TGT-P135-005}
Now \NSi{NS-F-P135-023}{b_s=2k'\eta/E}, so \eqref{eq:A.26} gives
\NSi{NS-F-P135-024}{|b_sw|\leq Ce_d}. Also, \NSi{NS-F-P135-025}{b_s^2\ll1} because
\NSi{NS-F-P135-026}{P_*>e^{T_d}}. Since \NSi{NS-F-P135-027}{a=2}, first choose
\NSi{NS-F-P135-028}{M_d} so that \NSi{NS-F-P135-029}{e_d} is small, then choose
\NSi{NS-F-P135-030}{P_*}. Both strict inequalities in \eqref{eq:A.24} follow.
\end{EESegment}

\begin{EESegment}{EE-TGT-P135-006}
\emph{The intermediate power-law interval.} During the axial transition, write
\NSi{NS-F-P135-031}{\mathcal A_X(U)=\bar k\eta}. Then
\NSi{NS-F-P135-032}{\bar k'=k-\bar k}. On the final eleven units, where \NSi{NS-F-P135-033}{k=0}, this gives
\NSi{NS-F-P135-034}{L\bar k<1/2}. It follows that \NSi{NS-F-P135-035}{W=1-L\bar k\geq1/2} on the next
transition. Put \NSi{NS-F-P135-036}{c_\eta=D\eta J_0'\asymp\eta^2}. Since \NSi{NS-F-P135-037}{U=0}, the angular
source on that transition is \NSi{NS-F-P135-038}{(-l)W-h+c_\eta}. The function
\NSi{NS-F-P135-039}{Q_s} keeps a positive lower bound that does not depend on \NSi{NS-F-P135-040}{\lambda}. The bound
for \NSi{NS-F-P135-041}{w} depends only on parameters chosen before \NSi{NS-F-P135-042}{\lambda}. So
\NSi{NS-F-P135-043}{\sqrt\lambda|w|} is small, while \NSi{NS-F-P135-044}{b_s=0}.
\end{EESegment}

\begin{EESegment}{EE-TGT-P135-007}
On the constant-slope interval, use its local logarithmic radial coordinate. There
\NSi{NS-F-P135-045}{\bar k(y)=\bar k(0)e^{-y}}, so the angular source is exactly
\NSi{NS-F-P135-046}{\lambda-h+c_\eta-\lambda L\bar k(0)e^{-y}}. Integrating the scalar equation for
\NSi{NS-F-P135-047}{Q_s} gives
\end{EESegment}
\NSFormula{NS-F-P135-048}{display}{A.27}
\begin{equation}
\begin{aligned}
 Q_s&\geq c_{\mathrm{pre}}(\eta^2+\lambda+e^{-(1-\lambda)y}),\\
 Q_s-\frac{\lambda-h+c_\eta}{1-\lambda}
   &=O_{C^1_\eta}\!\left(C_{\mathrm{pre}}(1+y)e^{-(1-\lambda)y}\right),\\
 |N_s/E|&\leq C_{\mathrm{pre}}|\eta|(1+y)e^{-(1/2-\lambda)y}.
\end{aligned}
\tag{A.27}\label{eq:A.27}
\end{equation}
\begin{EESegment}{EE-TGT-P135-008}
For the last line, \NSi{NS-F-P135-049}{U=0} makes the geometric source vanish, while the pressure source remains
\NSi{NS-F-P135-050}{O(E^2|\eta|)}. Convolving it with \NSi{NS-F-P135-051}{e^{-y}} gives the stated bound. The same
estimate holds after one \NSi{NS-F-P135-052}{\eta} derivative, but without the factor
\NSi{NS-F-P135-053}{|\eta|}. Using
\NSi{NS-F-P135-054}{|\eta|/(\eta^2+e^{-(1-\lambda)y})\leq\tfrac12e^{(1-\lambda)y/2}} now gives
\end{EESegment}
\NSFormula{NS-F-P135-055}{display}{none}
\[
 \sqrt\lambda|w|\leq C_{\mathrm{pre}}\sqrt\lambda(1+y)e^{\lambda y/2}=o(1)
 \qquad(0\leq y\leq60\log(1/\lambda)).
\]
\begin{EESegment}{EE-TGT-P135-009}
Since \NSi{NS-F-P135-056}{a=2+2\lambda} and \NSi{NS-F-P135-057}{b_s=0}, the stress lies in the admissible cone on
this interval. The same argument works on the preceding transition wherever \NSi{NS-F-P135-058}{l<0}; at its first
endpoint it gives the relaxed condition.
\end{EESegment}

\begin{EESegment}{EE-TGT-P135-010}
\emph{The pulse.} Both shear components change here, so the cone test needs a more exact relation between the averaged
axial velocity and the pulse profile. Set \NSi{NS-F-P135-059}{m=\mathcal A_X(U)/E} and
\NSi{NS-F-P135-060}{\beta=1/2-\lambda}. The exact equation for \NSi{NS-F-P135-061}{m} is
\NSi{NS-F-P135-062}{m'+\beta m=R_b}. Extend the main pulse by zero to the left and expand its exponential convolution
twice. The result is
\end{EESegment}
\NSFormula{NS-F-P135-063}{display}{A.28}
\begin{equation}
 m=\frac{R_b}{\beta}-\frac{\lambda\partial_{\zeta_b}R_b}{\beta^2}
   +O(\lambda^2)+m(0)e^{-\beta y}.
 \tag{A.28}\label{eq:A.28}
\end{equation}
\NSPage{136}%
\begin{EESegment}{EE-TGT-P136-001}
Taylor's formula gives this expansion because it bounds the convolution remainder by
\NSi{NS-F-P136-001}{C\lambda^2\lVert\partial_{\zeta_b}^2R_b\rVert_\infty
\int_0^\infty e^{-\beta v}v^2\,dv}. The pulse starts flat, and the end bumps have fixed
\NSi{NS-F-P136-002}{\zeta_b} derivatives that are exponentially small. The same expansion remains true after one
\NSi{NS-F-P136-003}{\eta} derivative. Equation~\eqref{eq:A.14} bounds the initial value.
\end{EESegment}

\begin{EESegment}{EE-TGT-P136-002}
In the angular equation, \NSi{NS-F-P136-004}{W=1-2D\eta Em-d(Em)_\eta}. Together with \eqref{eq:A.20}, this gives
\NSi{NS-F-P136-005}{S_q=\lambda-h+c_\eta+O(C_{\mathrm{pre}}E)}. Equations~\eqref{eq:A.14} and
\eqref{eq:A.27} then give
\end{EESegment}
\NSFormula{NS-F-P136-006}{display}{A.29}
\begin{equation}
 Q_s=\frac{\lambda-h+c_\eta}{1-\lambda}+O(C_{\mathrm{pre}}\lambda^{28}).
 \tag{A.29}\label{eq:A.29}
\end{equation}
\begin{EESegment}{EE-TGT-P136-003}
The pressure is normalized to vanish at radial infinity. It and its first
\NSi{NS-F-P136-007}{\eta} derivative are both \NSi{NS-F-P136-008}{O(E^2)}. Also, integrate
\NSi{NS-F-P136-009}{XE^2(R_b^2-1/2)} together with the difference of moments already built up before the pulse. This
shows that
\NSi{NS-F-P136-010}{S/(XE^2)=O_{C^1_\eta}(C_{\mathrm{pre}}e^{26}/\lambda)} throughout the pulse.
Formula~\eqref{eq:4.16} writes \NSi{NS-F-P136-011}{Q_s,N_s} in terms of the radial moment functions and their
\NSi{NS-F-P136-012}{\eta}-derivatives; it does not use radial derivatives of \NSi{NS-F-P136-013}{E,U}. Substituting
the estimates into that formula gives
\end{EESegment}
\NSFormula{NS-F-P136-014}{display}{A.30}
\begin{equation}
 \frac{N_s}{E}=-R_b+(D+c_\eta)m-D\eta m_\eta+O(C_{\mathrm{pre}}\lambda^{27}).
 \tag{A.30}\label{eq:A.30}
\end{equation}
\begin{EESegment}{EE-TGT-P136-004}
Every constant here is independent of \NSi{NS-F-P136-015}{\mathsf C}. Since
\NSi{NS-F-P136-016}{|\eta|/(\lambda+c_\eta)\leq C/\sqrt\lambda}, equation~\eqref{eq:A.20} makes the term from
\NSi{NS-F-P136-017}{\eta m_\eta/Q_s} an \NSi{NS-F-P136-018}{o(1)} contribution. Combining
\eqref{eq:A.28}--\eqref{eq:A.30} gives
\end{EESegment}
\NSFormula{NS-F-P136-019}{display}{none}
\[
 w=2R_b-C_d\partial_{\zeta_b}R_b+o(1),
 \qquad C_d=\frac{\lambda(D+c_\eta)(1-\lambda)}{\beta^2(\lambda-h+c_\eta)}\leq2+o(1).
\]
\begin{EESegment}{EE-TGT-P136-005}
This estimate for \NSi{NS-F-P136-020}{w} turns the cone test into bounds on the pulse shape and its derivative. For the
coefficient of \NSi{NS-F-P136-021}{R_b}, use
\NSi{NS-F-P136-022}{D+c_\eta-\beta=\lambda-h+c_\eta}. To bound
\NSi{NS-F-P136-023}{C_d}, maximize the quotient at \NSi{NS-F-P136-024}{c_\eta=0} and use
\NSi{NS-F-P136-025}{h/\lambda\to0}. Also, \NSi{NS-F-P136-026}{a=2+2\lambda} and
\NSi{NS-F-P136-027}{b_s=-R_b+O(\lambda)}. Apart from the exponentially small end bumps,
\NSi{NS-F-P136-028}{R_b\geq0} and \NSi{NS-F-P136-029}{\partial_{\zeta_b}R_b\leq1.2}, so
\end{EESegment}
\NSFormula{NS-F-P136-030}{display}{none}
\[
\begin{aligned}
 b_sw&\leq\sup_{R\geq0}(-2R^2+2.41R)+o(1)<.74,\\
 2b_sw+b_s^2/a+(a-2)w^2
   &\leq\sup_{R\geq0}(-3.5R^2+4.82R)+o(1)<1.68.
\end{aligned}
\]
\begin{EESegment}{EE-TGT-P136-006}
These quadratic cross terms use only the upper bound on the pulse derivative. Near the cutoff the derivative is negative,
which makes the left sides smaller. Every derivative remains bounded by a fixed constant. The first margin gives
\NSi{NS-F-P136-031}{a-b_sw>1}, while \NSi{NS-F-P136-032}{v_s\geq a>2}. This proves that the stress stays in the
admissible cone throughout the pulse.
\end{EESegment}

\begin{EESegment}{EE-TGT-P136-007}
\emph{Profile interpolation and exterior transition.} After the pulse,
\NSi{NS-F-P136-033}{U=M=J=0}, so \NSi{NS-F-P136-034}{W=1}. During the profile interpolation, the angular source is
\NSi{NS-F-P136-035}{-l-h+\vartheta_f c_\eta}. It follows that \NSi{NS-F-P136-036}{Q_s\geq c\lambda}, and the small
\NSi{NS-F-P136-037}{C^1_\eta} angular correction keeps that bound. Since
\NSi{NS-F-P136-038}{S(\infty)=0},
\end{EESegment}
\NSFormula{NS-F-P136-039}{display}{A.31}
\begin{equation}
 \frac{S(X)}{X}=\frac1{2X}\int_X^\infty E(X')^2\,dX'.
 \tag{A.31}\label{eq:A.31}
\end{equation}
\begin{EESegment}{EE-TGT-P136-008}
Until the exterior transition begins, \NSi{NS-F-P136-040}{l\leq-\lambda/2}; the remaining integral over the exterior
transition is bounded by \eqref{eq:A.18}. This gives
\NSi{NS-F-P136-041}{(|S|+|S_\eta|)/X\leq CE^2/\lambda}, including the angular edit, while
\NSi{NS-F-P136-042}{|\Pi|+|\Pi_\eta|\leq CE^2}. Formula~\eqref{eq:4.16} now gives
\NSi{NS-F-P136-043}{|w|\leq CE/\lambda^2\ll1}, because the end of the pulse made
\NSi{NS-F-P136-044}{E} exponentially small in \NSi{NS-F-P136-045}{1/\lambda}. Here
\NSi{NS-F-P136-046}{b_s=0} and \NSi{NS-F-P136-047}{2<a\leq2+2\lambda+.2+o(1)}, so both cone margins remain
positive.
\end{EESegment}

\begin{EESegment}{EE-TGT-P136-009}
Once the exterior transition has begun, its fields and the remaining radial moment functions are uniform in
\NSi{NS-F-P136-048}{\eta}. The only exceptions are the explicit coefficients in the formula for
\NSi{NS-F-P136-049}{Q_s,N_s}. Equations~\eqref{eq:A.23} and \eqref{eq:A.31} together give the following exact
expression.
\end{EESegment}
\NSPage{137}%
\NSFormula{NS-F-P137-001}{display}{none}
\[
 N_s=2\eta\int_y^\infty\bigl(he^{v-y}-A\bigr)E(v)^2\,dv=O(E(y)^2).
\]
\begin{EESegment}{EE-TGT-P137-001}
Throughout the exterior transition and at every larger radius,
\NSi{NS-F-P137-002}{l\leq-3h/4}. Since \NSi{NS-F-P137-003}{(\log E)'=l-1/2}, for
\NSi{NS-F-P137-004}{v\geq y} this gives
\end{EESegment}
\NSFormula{NS-F-P137-005}{display}{none}
\[
 E(v)^2=E(y)^2\exp\!\left(\int_y^v(2l(s)-1)\,ds\right)
 \leq E(y)^2e^{-(1+3h/2)(v-y)}.
\]
\begin{EESegment}{EE-TGT-P137-002}
It follows that
\end{EESegment}
\NSFormula{NS-F-P137-006}{display}{none}
\[
\begin{aligned}
 h\int_y^\infty e^{v-y}E(v)^2\,dv
   &\leq hE(y)^2\int_0^\infty e^{-3hs/2}\,ds=\frac23E(y)^2,\\
 \int_y^\infty E(v)^2\,dv
   &\leq E(y)^2\int_0^\infty e^{-(1+3h/2)s}\,ds
     =\frac{E(y)^2}{1+3h/2}.
\end{aligned}
\]
\begin{EESegment}{EE-TGT-P137-003}
Together with \NSi{NS-F-P137-007}{A=1/2+h} and \NSi{NS-F-P137-008}{|\eta|\leq1}, these estimates prove the stated
bound for \NSi{NS-F-P137-009}{N_s}, uniformly in \NSi{NS-F-P137-010}{h}. Before the steep-decay interval,
\NSi{NS-F-P137-011}{Q_s\geq c\lambda}. After that interval, and through the first half-unit of the terminal interval,
\NSi{NS-F-P137-012}{Q_s\geq ch}. On the last constant-slope interval it falls only to
\NSi{NS-F-P137-013}{Q_p\asymp h}, and \eqref{eq:A.16} gives the same lower bound for
\NSi{NS-F-P137-014}{0\leq y\leq1/2}. Equation~\eqref{eq:A.17} gives
\NSi{NS-F-P137-015}{E\leq Ce_{\mathrm{rel}}h^6} after the steep-decay interval. So
\NSi{NS-F-P137-016}{E/Q_s\ll1}, even when \NSi{NS-F-P137-017}{h} is much smaller than
\NSi{NS-F-P137-018}{e_{\mathrm{rel}}}. It follows that \NSi{NS-F-P137-019}{|w|\ll1},
\NSi{NS-F-P137-020}{b_s=0}, and \NSi{NS-F-P137-021}{2<a\leq4} on the part of the exterior transition being
considered. This proves the strict true cone there.
\end{EESegment}

\begin{EESegment}{EE-TGT-P137-004}
Every lower bound for \NSi{NS-F-P137-022}{Q_s} used above holds on a fixed compact logarithmic interval after all the
choices in \eqref{eq:A.6}. On such an interval, the fields and the functions
\NSi{NS-F-P137-023}{Q_s,N_s} do not depend on \NSi{NS-F-P137-024}{X_R} under the normalizations
\eqref{eq:A.4}. One final increase of \NSi{NS-F-P137-025}{X_R} makes the test that needs large
\NSi{NS-F-P137-026}{p_{s,1}} hold uniformly. This completes the proof of Proposition~\ref{prop:A.4}.
\end{EESegment}

\begin{EESegment}{EE-TGT-P137-005}
\subsection{Replacing the exterior power law with a heat flow.}\label{sec:A.6}
\end{EESegment}
\begin{EESegment}{EE-TGT-P137-006}
The reference profile \NSi{NS-F-P137-027}{(U,E)} from Proposition~\ref{prop:A.4} has the required moments. Past the
terminal interval it also has \NSi{NS-F-P137-028}{U=0} and
\NSi{NS-F-P137-029}{E=E_{\mathrm{pow}}=c_\infty X^{-A}}. This exterior power law
\NSi{NS-F-P137-030}{E_{\mathrm{pow}}} still leaves a viscous residual. The next step is to replace it outside the annulus
with the exact radial swirl heat flow \NSi{NS-F-P137-031}{K} from Lemma~\ref{lem:A.6}. After the replacement, the
exterior momentum residual is zero, and away from the axis the field has smooth limits at the singular time. The
replacement changes three moments. Proposition~\ref{prop:A.7} restores them with a correction in the second reserved
interval \NSi{NS-F-P137-032}{\mathcal I_2}; this also restores the pressure datum
\NSi{NS-F-P137-033}{\Pi_0} already prescribed for the axis in \eqref{eq:A.21}.
\end{EESegment}

\begin{EESegment}{EE-TGT-P137-007}
For a profile that is regular at the axis, let
\NSi{NS-F-P137-034}{\mathscr R_\theta^{(0)},\mathscr R_z^{(0)}} be its leading azimuthal and axial momentum residuals.
In the notation of \eqref{eq:4.12}, this means \NSi{NS-F-P137-035}{\mathscr R_j^{(0)}=R_j^{(0)}}. These residuals
include radial viscosity but leave out the higher-order axial viscosity. As in Proposition~\ref{prop:4.2}, the stress
components are \NSi{NS-F-P137-036}{T_\theta(r)=-r^{-2}\int_0^r s^2\mathscr R_\theta^{(0)}(s)\,ds} and
\NSi{NS-F-P137-037}{T_z(r)=-r^{-1}\int_0^r s\mathscr R_z^{(0)}(s)\,ds}. The next subsection uses the corrected moment
conditions to turn these into integrals from \NSi{NS-F-P137-038}{r} to infinity, in Lemma~\ref{lem:A.8}. It then finds
the sign and limiting direction of the stress as the outer endpoint is approached, in Proposition~\ref{prop:A.10}.
\end{EESegment}

\begin{EESegment}{EE-TGT-P137-008}
Throughout this subsection, \NSi{NS-F-P137-039}{0<h<1/2}, \NSi{NS-F-P137-040}{A=1/2+h}, and
\NSi{NS-F-P137-041}{D=1/2-h}. Keep \NSi{NS-F-P137-042}{d=1-\eta^2},
\NSi{NS-F-P137-043}{\tau=1-t}, and \NSi{NS-F-P137-044}{L=1-2h\eta^2} from \eqref{eq:4.1}. Every parameter in the
outer radial profile is fixed in the order shown in \eqref{eq:A.6}, before \NSi{NS-F-P137-045}{X_R} is increased.
On the tail from \eqref{eq:A.12}, use the local logarithmic radial coordinate
\NSi{NS-F-P137-046}{y=\log(X/X_{\mathrm{tail}})}. Write
\NSi{NS-F-P137-047}{X_b=e^3X_{\mathrm{tail}}} and
\NSi{NS-F-P137-048}{E_{\mathrm{pow}}=c_\infty X^{-A}}, where
\NSi{NS-F-P137-049}{c_\infty>0} does not depend on \NSi{NS-F-P137-050}{\eta}. The reference tail is then
\NSi{NS-F-P137-051}{E_{\mathrm{pow}}f_o(y)}.
\end{EESegment}

\NSPage{138}%
\begin{EESegment}{EE-TGT-P138-001}
Set \NSi{NS-F-P138-001}{a_K=1+h}. Define the heat factor by
\end{EESegment}
\NSFormula{NS-F-P138-002}{display}{A.32}
\begin{equation}
 \mathcal H(Z)=\frac1{\Gamma(a_K)}\int_0^\infty
 e^{-v}v^{a_K-1}(1+Zv)^{-h}\,dv,\qquad Z\geq0,
 \tag{A.32}\label{eq:A.32}
\end{equation}
\begin{EESegment}{EE-TGT-P138-002}
Here \NSi{NS-F-P138-003}{\Gamma(a_K)=\int_0^\infty e^{-v}v^{a_K-1}\,dv} is the gamma integral, so
\NSi{NS-F-P138-004}{\mathcal H(0)=1}. The heat profile is
\NSi{NS-F-P138-005}{E_{\mathrm{heat}}(X,\eta)=E_{\mathrm{pow}}(X)\mathcal H(2d/X)}.
\end{EESegment}

\begin{lemma}\label{lem:A.6}
\begin{EESegment}{EE-TGT-P138-003}
The function \NSi{NS-F-P138-006}{\mathcal H} is positive and is smooth from the right all the way to
\NSi{NS-F-P138-007}{Z=0}. Every fixed derivative stays bounded on each bounded nonnegative interval. The field
\end{EESegment}
\NSFormula{NS-F-P138-008}{display}{A.33}
\begin{equation}
 K(r,t)=c_\infty s^{-A}\mathcal H(2\tau/s),\qquad s=r^2/2,
 \tag{A.33}\label{eq:A.33}
\end{equation}
\begin{EESegment}{EE-TGT-P138-004}
does not depend on \NSi{NS-F-P138-009}{z}. It satisfies
\NSi{NS-F-P138-010}{\partial_tK=(\partial_{rr}+r^{-1}\partial_r-r^{-2})K}, and
\NSi{NS-F-P138-011}{K_r<0}. Its profile is
\NSi{NS-F-P138-012}{E_{\mathrm{pow}}(X)\mathcal H(2d/X)}. This profile is smooth, and all of its
\NSi{NS-F-P138-013}{\eta}-derivatives extend continuously from the interior to
\NSi{NS-F-P138-014}{\eta=\pm1}.
\end{EESegment}
\end{lemma}
\begin{proof}
\begin{EESegment}{EE-TGT-P138-005}
For the rising factorial \NSi{NS-F-P138-015}{(b)_m=b(b+1)\cdots(b+m-1)}, with
\NSi{NS-F-P138-016}{(b)_0=1}, differentiation under the integral, justified by a dominating integrable function,
gives
\end{EESegment}
\NSFormula{NS-F-P138-017}{display}{A.34}
\begin{equation}
 \mathcal H^{(m)}(Z)=\frac{(-1)^m(h)_m}{\Gamma(1+h)}
 \int_0^\infty e^{-v}v^{h+m}(1+Zv)^{-h-m}\,dv,
 \tag{A.34}\label{eq:A.34}
\end{equation}
\NSFormula{NS-F-P138-018}{display}{A.35}
\begin{equation}
 \mathcal H^{(m)}(0)=(-1)^m(h)_m(1+h)_m.
 \tag{A.35}\label{eq:A.35}
\end{equation}
\begin{EESegment}{EE-TGT-P138-006}
Deleting the last factor leaves an integrable function that dominates the integrand, including at
\NSi{NS-F-P138-019}{Z=0}. So the usual finite Taylor formula holds there for every fixed derivative order; a
convergent Taylor series is not needed. On any fixed small nonnegative interval, it gives
\end{EESegment}
\NSFormula{NS-F-P138-020}{display}{A.36}
\begin{equation}
 \mathcal H(Z)=1-h(1+h)Z+O_h(Z^2),\qquad
 |\mathcal H(Z)-1|+|Z\mathcal H'(Z)|\leq C_hZ.
 \tag{A.36}\label{eq:A.36}
\end{equation}
\begin{EESegment}{EE-TGT-P138-007}
The last constant can be chosen uniformly for \NSi{NS-F-P138-021}{0<h\leq h_0<1/2}. The reason is that
\NSi{NS-F-P138-022}{(h)_m(1+h)_m\leq C_mh} for every fixed \NSi{NS-F-P138-023}{m\geq1}. Now integrate
\NSi{NS-F-P138-024}{h\partial_v[e^{-v}v^{a_K}(1+Zv)^{-a_K}]}. This gives
\end{EESegment}
\NSFormula{NS-F-P138-025}{display}{A.37}
\begin{equation}
 Z^2\mathcal H''+(1+2a_KZ)\mathcal H'+a_K(a_K-1)\mathcal H=0;
 \tag{A.37}\label{eq:A.37}
\end{equation}
\begin{EESegment}{EE-TGT-P138-008}
both boundary terms are zero. Differentiate \eqref{eq:A.33} directly, using
\NSi{NS-F-P138-026}{Z=2\tau/s}. The two sides of the heat equation become
\end{EESegment}
\NSFormula{NS-F-P138-027}{display}{none}
\[
\begin{aligned}
 \partial_tK&=-2c_\infty s^{-A-1}\mathcal H',\\
 (\partial_{rr}+r^{-1}\partial_r-r^{-2})K
 &=2c_\infty s^{-A-1}
   [Z^2\mathcal H''+(2A+1)Z\mathcal H'+(A^2-1/4)\mathcal H].
\end{aligned}
\]
\begin{EESegment}{EE-TGT-P138-009}
Since \NSi{NS-F-P138-028}{2A+1=2a_K} and
\NSi{NS-F-P138-029}{A^2-1/4=a_K(a_K-1)}, equation~\eqref{eq:A.37} proves the heat equation, including its swirl
term. Next, \NSi{NS-F-P138-030}{-Z\mathcal H'/\mathcal H} is the average of
\NSi{NS-F-P138-031}{hZv/(1+Zv)} with respect to the positive probability density proportional to the integrand in
\eqref{eq:A.32}. So
\end{EESegment}
\NSFormula{NS-F-P138-032}{display}{none}
\[
 0\leq-\frac{Z\mathcal H'}{\mathcal H}<h,
 \qquad
 \frac{rK_r}{2K}=-A-\frac{Z\mathcal H'}{\mathcal H}<-\frac12.
\]
\begin{EESegment}{EE-TGT-P138-010}
The composition with \NSi{NS-F-P138-033}{Z=2(1-\eta^2)/X} only uses
\NSi{NS-F-P138-034}{Z\geq0}. Every nonzero \NSi{NS-F-P138-035}{\eta}-derivative is a finite sum of terms
\NSi{NS-F-P138-036}{\mathcal H^{(k)}(2d/X)(2/X)^k}, multiplied by polynomials in
\NSi{NS-F-P138-037}{d'=-2\eta} and \NSi{NS-F-P138-038}{d''=-2}. No term divides by
\NSi{NS-F-P138-039}{d}. This proves smoothness on the closed interval. The integral formula is used only on its
stated nonnegative domain. In particular, the large-\NSi{NS-F-P138-040}{X} expansion, including its remainder
after any fixed application of \NSi{NS-F-P138-041}{((X\partial_X)^j\partial_\eta^m)}, is
\end{EESegment}
\NSFormula{NS-F-P138-042}{display}{A.38}
\begin{equation}
 \frac{E_{\mathrm{heat}}(X,\eta)}{E_{\mathrm{pow}}(X)}
 =\mathcal H(2d/X)=1-\frac{2h(1+h)d}{X}+O_{h,j,m}(X^{-2}).
 \tag{A.38}\label{eq:A.38}
\end{equation}
\end{proof}
